\documentclass[a4paper,10pt]{report}

\def\packagepath{../../preambule}   % path du package princial
\usepackage{\packagepath/preambule}    % utilisation du fichier de configuration

\def\level{Première }              % Classe
\def\course{NSI}              % Matière
\def\eval{DS02 - Fonctions - IF}

\def\visibleornot{visible}    % visible or invisible
\def\documentpath{./Sources_Latex}
\def\date{Mardi 30 Septembre 2025}
\renewcommand{\arraystretch}{1}  % Ecart dans les tableau

\def\ptsexoA{4}
\def\ptsexoB{2}
\def\ptsexoC{2}
\def\ptsexoD{3}
\def\ptsexoE{2}
\def\ptsexoF{1}
\def\ptsexoG{2}

\def\ptstotal{\ptsexoA+\ptsexoB}

\usepackage{hyperref}

\usetikzlibrary{shapes, arrows, positioning}
\usetikzlibrary{shapes.geometric, arrows, positioning, decorations.pathreplacing}

\tikzstyle{etat} = [draw, rounded corners=15pt, minimum width=2.5cm, minimum height=1.2cm, text centered, font=\bfseries]
\tikzstyle{nouveau} = [etat, fill=blue!40]
\tikzstyle{pret} = [etat, fill=green!60]
\tikzstyle{elu} = [etat, fill=yellow!60]
\tikzstyle{bloque} = [etat, fill=red!60]
\tikzstyle{termine} = [etat, fill=white]
\tikzstyle{fleche} = [thick, ->, >=stealth]

\usepackage{float} % Permet d'utiliser [H]
\usepackage[T1]{fontenc}
\begin{document}

%%%%%%%%%%%%%%%%%%%%%%%%%%%%%%%%%%%%%%%%%%%%%%%%%%ù

%\PageGardeSujetBac[Matiere = Numérique et Sciences Informatiques,
%                  Session = 2025,
%                  AffJour=false,
%                  Duree = {3 heures 30},
%                  DernierePage = \pageref{LastPage},
%                  ModeExamen=false
%                  ]
\renewcommand{\labelitemi}{\textbullet} %pour éviter les tirets dans les "itemize" qui apportent confusion avec le signe moins.
%% Début page de garde style BAC

   %%%%%%%%%%%%%%%%%%%%%%%%%%%%%%%%%%%%%%%%%%%%%%%%%%%%%%%%
%\PageGardeSujetBac[clés]

\pagestyle{DS_FP}          % Feuille de style fancy
\NomPrenomNote{}


\exods{\ptsexoA}


On considère la fonction \verb|ma_fonction| suivante.

\begin{CodePythontexAlt}[TaillePolice=\large,Largeur=1\linewidth,Centre,PremLigne=1]{}
def ma_fonction(x, y):
    ''' fonction renvoyant la somme de x et y si x est plus grand que y 
    renvoie la différence sinon '''
    if x > y:
        sortie = x + y
    else:
        sortie = x - y 
    return sortie
\end{CodePythontexAlt}

\begin{enumerate}
    \item Quels sont les paramètres de la fonction \verb|ma_fonction|? \dotfill \verb|x| et \verb|y|
    \item Si je veux appliquer ma fonction aux valeurs 5 et 4. que dois-je taper dans la console Python? \dotfill \verb|ma_fonction(5, 4)|
    \item Comment se nomme la partie comprise entre les \verb|'''|? \dotfill \verb|Le Docstring|
    \item Remplir le tableau suivant lors de l'appel de la fonction \verb|ma_fonction(8, 5)|
    
    \begin{minipage}{0.6\linewidth}
        \begin{itemize}
        \item Dans la colonne ligne, indiquer à quelle ligne on se trouve
        \item Dans la colonne IF indiquer True ou False suivant si la ligne 4 est vraie ou fausse;
        \item Dans la colonne sortie, indiquer la valeur de la variable sortie.
    \end{itemize}
    \end{minipage}
    \begin{minipage}{0.39\linewidth}
         \begin{tabular}{|p{1.5cm}|p{1cm}|p{1.5cm}|}
        \hline
    Ligne & IF & sortie \\ \hline 
    4&True& \\ \hline       
    5&& 13\\ \hline       
    8&& 13\\ \hline       
    && \\ \hline       
    && \\ \hline       
    \end{tabular}
    \end{minipage}

    
   
\end{enumerate}


\bigskip

\exods{\ptsexoB}

Ecrire le code de la fonction \verb|moyenne(a, b, c)| qui renvoie la moyenne des valeurs $a$, $b$ et $c$.


\begin{CodePythontexAlt}[TaillePolice=\large,Largeur=1\linewidth,Centre,PremLigne=1]{}
def moyenne(a, b, c):
    ''' Renvoie la moyenne de a, b et c '''
    somme_valeurs = a + b + c
    return somme_valeurs / 4
\end{CodePythontexAlt}


\bigskip

\exods{\ptsexoC}

Ecrire le code de la fonction \verb|est_plus_long(mot1, mot2)| qui renvoie \verb|True| si \verb|mot1| est plus long que \verb|mot2| et False sinon.

On supposera que \verb|mot1| et \verb|mot2| sont deux chaines de caractères. On utilisera obligatoirement un IF.

\begin{CodePythontexAlt}[TaillePolice=\large,Largeur=1\linewidth,Centre,PremLigne=1]{}
def est_plus_long(mot1, mot2):
    ''' Renvoie True si mot1 et plus long que mot2 et False sinon '''
    longueur_mot1 = len(mot1)
    longueur_mot2 = len(mot2)
    if longueur_mot1 > longueur_mot2:
        return True
    else:
        return False
\end{CodePythontexAlt}


\bigskip



\end{document}








